\section{Design Context and Drivers}

\subsection{Problem Statement}
This report documents the engineering reasoning behind a proposed design. It is
best used when reviewers need to understand requirements, interfaces,
constraints, and why a specific architecture was selected over alternatives.

\specbox{
Design reports should explain \emph{why this approach exists} before they dive
into schematics, algorithms, or dimensions. A clear problem statement reduces
review churn more than an extra page of equations.
}

\subsection{Primary Drivers}
\begin{table}[H]
    \centering
    \small
    \begin{tabular}{l p{9.5cm}}
        \toprule
        \textbf{Driver} & \textbf{Interpretation} \\
        \midrule
        Functional Objective & What the design must do at the system level. \\
        Performance Constraint & Accuracy, speed, efficiency, timing, or capacity requirements. \\
        Physical Constraint & Size, weight, thermal, EMC, or environmental limits. \\
        Cost / Supply Constraint & Approved part families, sourcing restrictions, or target cost. \\
        Integration Constraint & Interfaces to firmware, mechanics, manufacturing, or test. \\
        \bottomrule
    \end{tabular}
    \caption{Recommended design-driver summary}
    \label{tab:design-drivers}
\end{table}

\subsection{Assumptions}
\begin{itemize}
    \item state any inherited platform, firmware, or supplier assumption;
    \item call out which interfaces are fixed and which may still change;
    \item note any placeholder component, model, or preliminary geometry.
\end{itemize}
